\documentclass[
	paper=128mm:96mm,
	fontsize=11pt,
	pagesize,
	parskip=half-,
]{scrartcl}

\linespread{1}

%------------------------------------------------
% Colors
\usepackage{xcolor}	 % Required for custom colors
% Define a few colors for making text stand out within the presentation
\definecolor{mygreen}{RGB}{44,85,17}
\definecolor{myblue}{RGB}{34,31,217}
\definecolor{mybrown}{RGB}{194,164,113}
\definecolor{myred}{RGB}{255,66,56}
% Use these colors within the presentation by enclosing text in the commands below
\newcommand*{\mygreen}[1]{\textcolor{mygreen}{#1}}
\newcommand*{\myblue}[1]{\textcolor{myblue}{#1}}
\newcommand*{\mybrown}[1]{\textcolor{mybrown}{#1}}
\newcommand*{\myred}[1]{\textcolor{myred}{#1}}
%------------------------------------------------

%------------------------------------------------
% Margins
\usepackage[ % Page margins settings
	includeheadfoot,
	top=3.5mm,
	bottom=3.5mm,
	left=5.5mm,
	right=5.5mm,
	headsep=6.5mm,
	footskip=8.5mm
]{geometry}
%------------------------------------------------

%------------------------------------------------
% Fonts
\usepackage[T1]{fontenc}
\usepackage{lmodern}
\renewcommand{\familydefault}{\sfdefault}
%------------------------------------------------

%------------------------------------------------
% Various required packages
\usepackage{amsthm}
\usepackage{bm}
\usepackage{graphicx}
\usepackage{tikz}
\usepackage{booktabs}
\usepackage{multicol}
\usepackage{lastpage}
\usepackage[english]{babel}
\usepackage{microtype}
\usepackage{tocstyle}
\usepackage[backend=biber,sorting=none,natbib=true,style=numeric]{biblatex}
\usepackage{wallpaper}
\usepackage{epstopdf}
\usepackage{amssymb, amsmath}
\usepackage{pdfpages}
\usepackage[english]{babel}
\usepackage{empheq}
\usepackage{multicol}
\usepackage{multirow}
\usepackage{cases}
\usepackage{scrextend}
\usepackage{amsfonts}
\usepackage{amsmath}
\usepackage{listings}
\usepackage{tikz}
\usepackage{float}
\usetikzlibrary{shapes,arrows,positioning,shadows,calc}
%------------------------------------------------

%------------------------------------------------
% Load bibliography
\bibliography{bib}
\addbibresource{bib}
%------------------------------------------------


%------------------------------------------------
% Slide layout configuration
\usepackage{scrpage2} % Required for customization of the header and footer
\pagestyle{scrheadings} % Activates the pagestyle from scrpage2 for custom headers and footers
\clearscrheadfoot % Remove the default header and footer
\setkomafont{pageheadfoot}{\normalfont\color{black}\sffamily} % Font settings for the header and footer

% Sets vertical centering of slide contents with increased space between paragraphs/lists
\makeatletter
\renewcommand*{\@textbottom}{\vskip \z@ \@plus 1fil}
\newcommand*{\@texttop}{\vskip \z@ \@plus .5fil}
\addtolength{\parskip}{\z@\@plus .25fil}
\makeatother

% Remove page numbers and the dots leading to them from the outline slide
\makeatletter
\newtocstyle[noonewithdot]{nodotnopagenumber}{\settocfeature{pagenumberbox}{\@gobble}}
\makeatother
\usetocstyle{nodotnopagenumber}

\AtBeginDocument{\renewcaptionname{english}{\contentsname}{\Large Outline}} % Change the name of the table of contents
%------------------------------------------------

%------------------------------------------------
% Header configuration - if you don't want a header remove this block
\ohead{
\hspace{-2mm}
\begin{tikzpicture}[remember picture, overlay]
	\node at (-\paperwidth/2+3.2cm,0) {\includegraphics{data/banner.png}};
\end{tikzpicture}
\color{black}\runninghead} % Header text defined by the \runninghead command below and colored white for contrast
%------------------------------------------------

%------------------------------------------------
% Footer configuration
\newlength{\footheight}
\setlength{\footheight}{8mm} % Height of the footer
\addtokomafont{pagefoot}{\footnotesize} % Small font size for the footnote

\ifoot{% Left side
\hspace{-2mm}
\begin{tikzpicture}[remember picture,overlay]
\node [xshift=\paperwidth/2,yshift=\footheight] at (current page.south west)[rectangle,fill,inner sep=0pt,minimum width=\paperwidth,minimum height=3pt,top color=mygreen,bottom color=mygreen,opacity=0.7]{}; % Green bar
\end{tikzpicture}
\myauthor\ \raisebox{0.2mm}{$\bm{\vert}$}\ \myuni % Left side text
}

\ofoot[\pagemark/\pageref{LastPage}\hspace{-2mm}]{\pagemark/\pageref{LastPage}\hspace{-2mm}} % Right side
%------------------------------------------------

%------------------------------------------------
% Section spacing - deeper section titles are given less space due to lesser importance
\usepackage{titlesec} % Required for customizing section spacing
\titlespacing{\section}{0mm}{0mm}{0mm} % Lengths are: left, before, after
\titlespacing{\subsection}{0mm}{0mm}{-1mm} % Lengths are: left, before, after
\titlespacing{\subsubsection}{0mm}{0mm}{-2mm} % Lengths are: left, before, after
\setcounter{secnumdepth}{0} % How deep sections are numbered, set to no numbering by default - change to 1 for numbering sections, 2 for numbering sections and subsections, etc
%------------------------------------------------

%------------------------------------------------
% Theorem style
\newtheoremstyle{mythmstyle} % Defines a new theorem style used in this template
{0.5em} % Space above
{0.5em} % Space below
{} % Body font
{} % Indent amount
{\sffamily\bfseries} % Head font
{} % Punctuation after head
{\newline} % Space after head
{\thmname{#1}\ \thmnote{(#3)}} % Head spec
	
\theoremstyle{mythmstyle} % Change the default style of the theorem to the one defined above
\newtheorem{theorem}{Theorem}[section] % Label for theorems
\newtheorem{remark}[theorem]{Remark} % Label for remarks
\newtheorem{algorithm}[theorem]{Algorithm} % Label for algorithms
\makeatletter % Correct qed adjustment
%------------------------------------------------

%------------------------------------------------
% The code for the box which can be used to highlight an element of a slide (such as a theorem)
\newcommand*{\mybox}[2]{ % The box takes two arguments: width and content
\par\noindent
\begin{tikzpicture}[mynodestyle/.style={rectangle,draw=mygreen,thick,inner sep=2mm,text justified,top color=white,bottom color=white,above}]\node[mynodestyle,at={(0.5*#1+2mm+0.4pt,0)}]{ % Box formatting
\begin{minipage}[t]{#1}
#2
\end{minipage}
};
\end{tikzpicture}
\par\vspace{-1.3em}}
%------------------------------------------------

%----------------------------------------------------------------------------------------
%	PRESENTATION INFORMATION
%----------------------------------------------------------------------------------------

\newcommand*{\mytitle}{Hive Simulation} % Title
\newcommand*{\runninghead}{Hive Simulation} % Running head displayed on almost all slides
\newcommand*{\myauthor}{Stefan Gugler, Elias Huwyler, Fabian Tschopp} % Presenters name(s)
\newcommand*{\mydate}{\today} % Presentation date
\newcommand*{\myuni}{ETH Zurich} % University or department

%----------------------------------------------------------------------------------------

\begin{document}

%----------------------------------------------------------------------------------------
%	TITLE SLIDE
%----------------------------------------------------------------------------------------

% Title slide - you may have to tweak a few of the numbers if you wish to make changes to the layout
\thispagestyle{empty} % No slide header and footer
\LRCornerWallPaper{1}{data/background.png}
\begin{tikzpicture}[remember picture,overlay] % Background box
	\node [xshift=\paperwidth/2,yshift=\paperheight/2] at (current page.south west)[rectangle,fill,inner sep=0pt,minimum width=\paperwidth,minimum height=\paperheight/3,top color=mygreen,bottom color=mygreen,opacity=0.7]{}; % Change the height of the box, its colors and position on the page here
	\node [xshift=1.7cm,yshift=-2.8cm] {\colorbox{white}{\includegraphics[scale=0.6]{data/ETHlogo.eps}}};
\end{tikzpicture}
% Text within the box
\begin{flushright}
	\vspace{0.6cm}
	\color{white}\sffamily
	{\bfseries\Large\mytitle\par}
	\vspace{0.5cm}
	\normalsize
	\myauthor\par
	\mydate\par
	\vfill
\end{flushright}

\clearpage
\ClearWallPaper

%----------------------------------------------------------------------------------------
%	TABLE OF CONTENTS
%----------------------------------------------------------------------------------------

\thispagestyle{empty}

\small\tableofcontents

\clearpage

%----------------------------------------------------------------------------------------
%	PRESENTATION SLIDES
%----------------------------------------------------------------------------------------
\section{Advanced model: Environment simulation}

	\clearpage

%------------------------------------------------

\subsection{Model overview}
	
% Define state styles
\tikzstyle{resource} = [circle, draw, top color=white, bottom color=red!20, draw=red!50!black!100, drop shadow, text width=1cm, text centered, node distance=2cm, minimum height=1cm]
\tikzstyle{bee} = [rectangle, draw, top color=white, bottom color=blue!30, draw=blue!50!black!100, drop shadow, text width=2cm, text centered, rounded corners, minimum height=1cm]
\tikzstyle{natural} = [rectangle, draw, top color=white, bottom color = green!30, draw=green!50!black!100, drop shadow, text width=2cm, text centered, rounded corners, minimum height=1cm]
\tikzstyle{line} = [line width=1.5pt, draw, -latex']


\begin{figure}[H]
	\centering
	\scalebox{0.6}{
	\large
	\begin{tikzpicture}[node distance = 3cm and 3cm, auto]
		% Place nodes
		\node [bee] (queen) {Queen};
		\node [bee, below of=queen, node distance = 2cm] (brood) {Uncapped brood};
		\node [bee, below of=brood, node distance = 2cm] (hiveBees) {Hive bees};
		\node [bee, below of=hiveBees, node distance = 2cm] (foragers) {Foragers and scouts};
		\node [resource, right of=brood, node distance = 5cm] (food) {Food};
		\node [natural, right of=food, node distance = 5cm] (environment) {Environment};
		\node [natural, below of=environment, node distance = 2cm] (flowers) {Flowers};
		\node [natural, above of=environment, node distance = 2cm] (seasons) {Seasons};
		% Draw edges
		\path [line] (queen.south) -- node {egg laying} (brood.north);
		\path [line] (brood.south) -- node {eclosion} (hiveBees.north);
		\path [line] ($(hiveBees.south)+(-0.25cm, 0.0cm)$) -- node [left] {food income requirements} ($(foragers.north)+(-0.25cm, 0.0cm)$);
		\path [line] ($(foragers.north)+(+0.25cm, 0.0cm)$) -- node [right] {social inhibition} ($(hiveBees.south)+(+0.25cm, 0.0cm)$);
		\path [line] (foragers.west) -- +(-1.0cm, 0.0cm) node [near end] {death};
		\path [line, dashed] (hiveBees.west) -- +(-0.5cm, 0.0cm) |-  node [near start] {affect survival} (brood.west);
		\path [line] (environment.north) -- node [left] {season durations}(seasons.south);
		\path [line] (environment.south) -- node [left] {patch distribution} (flowers.north);
		\path [line, dotted] (food) -- (queen.east);
		\path [line, dotted] (food) -- (hiveBees.east);
		\path [line, dotted] (food) -- (foragers.east);
		\path [line, dotted] (food) -- (brood.east);
		\path [line, dashed] (seasons.west) -- node {affect laying rate} (queen.east);
		\path [line, dashed] (seasons.east) -- +(0.5cm, 0.0cm) |- node [near start] {affect growth} (flowers.east);
		\path [line, dashed] (flowers.south) -- +(0.0cm, -2.5cm) node [above left] {foraging source availability} -| (foragers.south);
		\path [line, dotted] (foragers.east) -| node [near start, below] {food collection} (food.south);
	\end{tikzpicture}} 
	\caption{\textit{Honey bee social dynamics/influences covered by our model.}}
	\label{fig:modelOverview}
\end{figure}


	\clearpage

%------------------------------------------------

\subsection{Agents: Assigning jobs}
	% Define state styles
\tikzstyle{decision} = [diamond, draw, top color=white, bottom color=red!20, draw=red!50!black!100, drop shadow, text width=2cm, text centered, node distance=2cm, inner sep=0pt]
\tikzstyle{state} = [rectangle, draw, top color=white, bottom color=blue!30, draw=blue!50!black!100, drop shadow, text width=2cm, text centered, rounded corners, minimum height=2cm]
\tikzstyle{line} = [line width=1.5pt, draw, -latex']
\begin{figure}[H]
	\centering
	\scalebox{0.6}{
	\large
	\begin{tikzpicture}[node distance = 3cm and 3cm, auto]
		% Place nodes
		\node [state] (unemployed) {Unemployed bee};
		\node [decision, right of=unemployed, node distance = 3.5cm] (moreBees) {More active bees possible?};
		\node [decision, right of=moreBees, node distance = 4cm] (moreScouts) {More scouts possible?};
		\node [state, right of=moreScouts, node distance = 4cm] (foragerBee) {Forager bee, picking waggle dance randomly};
		\node [state, below of=moreScouts, node distance = 4cm] (randomWalk) {Scout bee, doing random walk};
		% Draw edges
		\path [line] +(-2cm, 0.0) -- (unemployed.west);
		\path [line] (unemployed) -- (moreBees);
		\path [line] (moreBees.north) -- +(0, +0.5cm) -| node [near start]{no} (unemployed.north);
		\path [line] (moreBees.east) -- node [near start]{yes} (moreScouts.west);
		\path [line] (moreScouts.east) -- node [near start] {no} (foragerBee.west);
		\path [line] (moreScouts.south) -- node [near start] {yes} (randomWalk.north);
	\end{tikzpicture}}
	\caption{\textit{Assigning jobs to unemployed bees. Scouts and foragers are possible.}}
\end{figure}


	\clearpage

%------------------------------------------------

\subsection{Foragers' distribution across flower patches}
	\begin{itemize}
		\item Bees will focus on newly reported and then on the most profitable food sources
		\item Driving factors for evaluating a patch ($p < q$):
		\begin{enumerate}
			\item Patch quality ($b_w$)
			\item Distance from the hive ($d_w$)
			\item Patch size ($A_w$)
		\end{enumerate}
		\item Relative evaluation based on what other foraging sources are available
	\end{itemize}
	\clearpage

%------------------------------------------------

\subsection{Agents: Scout bees}
	% Define state styles
\tikzstyle{decision} = [diamond, draw, top color=white, bottom color=red!20, draw=red!50!black!100, drop shadow, text width=2cm, text centered, node distance=2cm, inner sep=0pt]
\tikzstyle{state} = [rectangle, draw, top color=white, bottom color=blue!30, draw=blue!50!black!100, drop shadow, text width=2cm, text centered, rounded corners, minimum height=2cm]
\tikzstyle{line} = [line width=1.5pt, draw, -latex']
\begin{figure}[H]
	\centering
	\scalebox{0.6}{
	\large
	\begin{tikzpicture}[node distance = 3cm and 3cm, auto]
		% Place nodes
		\node [state, right of=moreScouts, node distance = 4cm] (randomWalk) {Scout bee, doing random walk};
		\node [decision, right of=randomWalk, node distance = 4cm] (foundPatch) {Flower patch found?};
		\node [decision, right of=foundPatch, node distance = 4cm] (limitReached) {Distance limit reached?};
		\node [state, right of=limitReached, node distance = 4cm] (scoutToHive) {Scout bee, returning to hive};
		% Draw edges
		\path [line] (randomWalk.east) -- (foundPatch.west);
		\path [line] (foundPatch.east) -- node [near start] {no} (limitReached.west);
		\path [line] (limitReached.east) -- node [near start] {yes} (scoutToHive.west);
		\path [line] (limitReached.north) -- node [near start] {no} +(0.0cm, 0.5cm) -| (randomWalk.north);
		\path [line] (foundPatch.south) -- +(0.0cm, -0.5cm) node [near start] {yes} -| (scoutToHive.south);
	\end{tikzpicture}}
	\caption{\textit{Scouting behaviour until a flower patch is found or the maximum distance is reached.}}
\end{figure}


	\clearpage


%------------------------------------------------

\subsection{Scouts' random walk}
	% Define state styles
\tikzstyle{decision} = [diamond, draw, top color=white, bottom color=red!20, draw=red!50!black!100, drop shadow, text width=2cm, text centered, node distance=2cm, inner sep=0pt]
\tikzstyle{state} = [rectangle, draw, top color=white, bottom color=blue!30, draw=blue!50!black!100, drop shadow, text width=2cm, text centered, rounded corners, minimum height=2cm]
\tikzstyle{line} = [line width=1.5pt, draw, -latex']
\begin{figure}[H]
	\centering
	\scalebox{0.6}{
	\large
	\begin{tikzpicture}[node distance = 3cm and 3cm, auto]
		% Place nodes
		\node [state, right of=moreScouts, node distance = 4cm] (randomWalk) {Scout bee, doing random walk};
		\node [decision, right of=randomWalk, node distance = 4cm] (foundPatch) {Flower patch found?};
		\node [decision, right of=foundPatch, node distance = 4cm] (limitReached) {Distance limit reached?};
		\node [state, right of=limitReached, node distance = 4cm] (scoutToHive) {Scout bee, returning to hive};
		% Draw edges
		\path [line] (randomWalk.east) -- (foundPatch.west);
		\path [line] (foundPatch.east) -- node [near start] {no} (limitReached.west);
		\path [line] (limitReached.east) -- node [near start] {yes} (scoutToHive.west);
		\path [line] (limitReached.north) -- node [near start] {no} +(0.0cm, 0.5cm) -| (randomWalk.north);
		\path [line] (foundPatch.south) -- +(0.0cm, -0.5cm) node [near start] {yes} -| (scoutToHive.south);
	\end{tikzpicture}}
	\caption{\textit{Scouting behaviour until a flower patch is found or the maximum distance is reached.}}
\end{figure}


	\clearpage


%------------------------------------------------

\subsection{Agents: Forager bees}
	% Define state styles
\tikzstyle{decision} = [diamond, draw, top color=white, bottom color=red!20, draw=red!50!black!100, drop shadow, text width=2cm, text centered, node distance=2cm, inner sep=0pt]
\tikzstyle{state} = [rectangle, draw, top color=white, bottom color=blue!30, draw=blue!50!black!100, drop shadow, text width=2cm, text centered, rounded corners, minimum height=2cm]
\tikzstyle{line} = [line width=1.5pt, draw, -latex']
\begin{figure}[H]
	\centering
	\scalebox{0.6}{
	\large
	\begin{tikzpicture}[node distance = 3cm and 3cm, auto]
		% Place nodes
		\node [state, right of=moreScouts, node distance = 4cm] (foragerBee) {Forager bee, picking waggle dance randomly};
		\node [decision, right of=foragerBee, node distance = 3cm] (pq1) {Evaluate $p < q$};
		\node [state, right of=pq1, node distance = 3cm] (foragerBeeToPatch) {Forager bee, on the way to the flower patch};
		\node [decision, right of=foragerBeeToPatch, node distance = 3.5cm] (foragerPatchReached) {Flower patch reached?};
		\node [state, right of=foragerPatchReached, node distance = 3.5cm] (collectingFood) {Forager bee, collecting food};
		\node [state, right of=collectingFood, node distance = 3.5cm] (foragerToHive) {Forager bee, returning to hive};
		% Draw edges
		\path [line] (foragerBee.east) -- (pq1.west);
		\path [line] (pq1.east) -- node [near start] {yes} (foragerBeeToPatch.west);
		\path [line] (pq1.north) -- node [near start] {no} +(0.0cm, +0.7cm) -| (foragerBee.north);
		\path [line] (foragerBeeToPatch.east) -- (foragerPatchReached.west);
		\path [line] (foragerPatchReached.east) -- node [near start] {yes} (collectingFood.west);
		\path [line] (foragerPatchReached.north) -- node [near start] {no} +(0.cm, +0.5cm) -| (foragerBeeToPatch.north);
		\path [line] (collectingFood.east) -- (foragerToHive.west);
	\end{tikzpicture}}
	\caption{\textit{Foraging behaviour.}}
\end{figure}


	\clearpage

%------------------------------------------------

\subsection{Path optimization}
	% Define state styles
\tikzstyle{decision} = [diamond, draw, top color=white, bottom color=red!20, draw=red!50!black!100, drop shadow, text width=2cm, text centered, node distance=2cm, inner sep=0pt]
\tikzstyle{state} = [rectangle, draw, top color=white, bottom color=blue!30, draw=blue!50!black!100, drop shadow, text width=2cm, text centered, rounded corners, minimum height=2cm]
\tikzstyle{line} = [line width=1.5pt, draw, -latex']
\begin{figure}[H]
	\centering
	\scalebox{0.6}{
	\large
	\begin{tikzpicture}[node distance = 3cm and 3cm, auto]
		% Place nodes
		\node [state, right of=moreScouts, node distance = 4cm] (randomWalk) {Scout bee, doing random walk};
		\node [decision, right of=randomWalk, node distance = 4cm] (foundPatch) {Flower patch found?};
		\node [decision, right of=foundPatch, node distance = 4cm] (limitReached) {Distance limit reached?};
		\node [state, right of=limitReached, node distance = 4cm] (scoutToHive) {Scout bee, returning to hive};
		% Draw edges
		\path [line] (randomWalk.east) -- (foundPatch.west);
		\path [line] (foundPatch.east) -- node [near start] {no} (limitReached.west);
		\path [line] (limitReached.east) -- node [near start] {yes} (scoutToHive.west);
		\path [line] (limitReached.north) -- node [near start] {no} +(0.0cm, 0.5cm) -| (randomWalk.north);
		\path [line] (foundPatch.south) -- +(0.0cm, -0.5cm) node [near start] {yes} -| (scoutToHive.south);
	\end{tikzpicture}}
	\caption{\textit{Scouting behaviour until a flower patch is found or the maximum distance is reached.}}
\end{figure}


	\clearpage

%------------------------------------------------

\subsection{Agents: Returning to the hive}
	
% Define state styles
\tikzstyle{decision} = [diamond, draw, top color=white, bottom color=red!20, draw=red!50!black!100, drop shadow, text width=2cm, text centered, node distance=2cm, inner sep=0pt]
\tikzstyle{state} = [rectangle, draw, top color=white, bottom color=blue!30, draw=blue!50!black!100, drop shadow, text width=2cm, text centered, rounded corners, minimum height=2cm]
\tikzstyle{line} = [line width=1.5pt, draw, -latex']
\begin{figure}
	\centering
	\scalebox{0.9}{
	\large
	\begin{tikzpicture}[node distance = 3cm and 3cm, auto]
		% Place nodes
		\node [state] (unemployed) {Unemployed bee};
		\node [decision, right of=unemployed, node distance = 3.5cm] (moreBees) {More active bees possible?};
		\node [decision, right of=moreBees, node distance = 4cm] (moreScouts) {More scouts possible?};
		\node [state, right of=moreScouts, node distance = 4cm] (foragerBee) {Forager bee, picking waggle dance randomly};
		\node [decision, below of=foragerBee, node distance = 3cm] (pq1) {Evaluate $p < q$};
		\node [state, below of=pq1, node distance = 3cm] (foragerBeeToPatch) {Forager bee, on the way to the flower patch};
		\node [decision, below of=foragerBeeToPatch, node distance = 3.5cm] (foragerPatchReached) {Flower patch reached?};
		\node [state, below of=foragerPatchReached, node distance = 3.5cm] (collectingFood) {Forager bee, collecting food};
		\node [state, below of=moreScouts, node distance = 3cm] (randomWalk) {Scout bee, doing random walk};
		\node [decision, below of=randomWalk, node distance = 3cm] (foundPatch) {Flower patch found?};
		\node [decision, below of=foundPatch, node distance = 3.5cm] (limitReached) {Distance limit reached?};
		\node [state, below of=limitReached, node distance = 3.5cm] (scoutToHive) {Scout bee, returning to hive};
		\node [state, below of=collectingFood, node distance = 3.5cm] (foragerToHive) {Forager bee, returning to hive};
		\node [decision, left of=scoutToHive, node distance = 4cm] (scoutHiveReached) {Hive reached?};
		\node [decision, left of=foragerToHive, node distance = 4cm] (foragerHiveReached) {Hive reached?};
		\node [decision, above of=scoutHiveReached, node distance = 6.5cm] (scoutingSuccess) {Scouting successful?};
		\node [state, above of=scoutingSuccess, node distance = 3.5cm] (scoutWaggle) {Scout bee, doing waggle dance};
		\node [state, left of=foragerHiveReached, node distance = 3cm] (foodUnloading) {Forager bee, food unloading};
		\node [decision, left of=foodUnloading, node distance = 3.5cm] (pq2) {Evaluate $p < q$};
		\node [decision, below of=unemployed, node distance = 9.5cm] (pq3) {Evaluate $p < q$};
		\node [state, below of=pq3, node distance = 3.5cm] (foragerWaggle) {Forager bee, doing waggle dance};
		% Draw edges
		\path [line] +(-2cm, 0.0) -- (unemployed.west);
		\path [line] (unemployed) -- (moreBees);
		\path [line] (moreBees.north) -- +(0, +0.5cm) -| node [near start]{no} (unemployed.north);
		\path [line] (moreBees.east) -- node [near start]{yes} (moreScouts.west);
		\path [line] (moreScouts.east) -- node [near start] {no} (foragerBee.west);
		\path [line] (foragerBee.south) -- (pq1.north);
		\path [line] (pq1.south) -- node [near start] {yes} (foragerBeeToPatch.north);
		\path [line] (pq1.east) -- +(+0.5cm, 0.0) |- node [near start] {no} (foragerBee.east);
		\path [line] (foragerBeeToPatch.south) -- (foragerPatchReached.north);
		\path [line] (foragerPatchReached.south) -- node [near start] {yes} (collectingFood.north);
		\path [line] (foragerPatchReached.east) -- node [near start] {no} +(+0.5cm, 0.0) |- (foragerBeeToPatch.east);
		\path [line] (moreScouts.south) -- node [near start] {yes} (randomWalk.north);
		\path [line] (randomWalk.south) -- (foundPatch.north);
		\path [line] (foundPatch.south) -- node [near start] {no} (limitReached.north);
		\path [line] (foundPatch.east) -- +(0.5cm, 0.0) |- node [near start] {yes} (scoutToHive.east);
		\path [line] (scoutToHive.west) -- (scoutHiveReached.east);
		\path [line] (collectingFood.south) -- (foragerToHive.north);
		\path [line] (foragerToHive.west) -- (foragerHiveReached.east);
		\path [line] (scoutHiveReached.south) -- +(0.0cm, -0.5cm) -| node [near start] {no} (scoutToHive.south);
		\path [line] (foragerHiveReached.south) -- +(0.0cm, -0.5cm) -| node [near start] {no} (foragerToHive.south);
		\path [line] (foragerHiveReached.west) -- node [near start] {yes} (foodUnloading);
		\path [line] (foodUnloading.west) -- (pq2);
		\path [line] (pq3.north) -- node [near start] {no} (unemployed);
		\path [line] (pq2.west) -- node [near start] {yes} +(-0.5cm, 0.0cm) |- +(0.0cm, 2.0cm) -| (foragerWaggle.south);
		\path [line] (foragerWaggle.north) -- (pq3.south);
		\path [line] (pq3.west) -- node [near start] {yes} +(-0.5cm, 0.0cm) |- +(0.0cm, -9cm) -| +(14.71cm, 0.0cm) |-  (foragerBeeToPatch);
		\path [line] (pq2.north) -- node [near start] {no} +(0.0cm, 0.65cm) -- +(1cm, 0.65cm) |- (pq3.east);
	\end{tikzpicture}}
	\caption{\textit{State transitions of the bees in the environmental simulation.}}
\end{figure}


	\clearpage

%------------------------------------------------

\subsection{Agent based model: recorded sample clips}
	\begin{itemize}
		\item Day 158, recorded sample with scouts displayed
		\item Day 158, recorded sample without scouts displayed
		\item Two different runs, not the same flower patches are being selected
	\end{itemize}
	\clearpage

%------------------------------------------------

\thispagestyle{empty}

\printbibliography[heading=none]

\clearpage

%------------------------------------------------

\thispagestyle{empty}
\LRCornerWallPaper{1}{data/background.png}

\begin{tikzpicture}[remember picture,overlay]
	\node [xshift=\paperwidth/2,yshift=\paperheight/2] at (current page.south west)[rectangle,fill,inner sep=0pt,minimum width=\paperwidth,minimum height=\paperheight/3,top color=mygreen,bottom color=mygreen,opacity=0.7]{};
\end{tikzpicture}

\begin{flushright}
	\vspace{0.6cm}
	\color{white}\sffamily
	{\bfseries\LARGE Questions?\par}
	\vfill
\end{flushright}

%----------------------------------------------------------------------------------------

\end{document}